\documentclass{article}
\usepackage{amsmath,amssymb,amsthm}
\usepackage{algorithm}
\usepackage[noend]{algpseudocode}
\usepackage{hyperref}
\usepackage{enumitem}
\usepackage{xcolor}
\newtheorem{theorem}{Theorem}
\newtheorem{lemma}{Lemma}
\newtheorem{assumption}{Assumption}
\newcommand{\E}{\mathbb{E}}
\newcommand{\R}{\mathbb{R}}

\begin{document}

\section*{RMSProp for Non‑Convex Smooth Optimization}

\section*{Mathematical Formulation}
Let $f:\R^{d}\rightarrow\R$ be differentiable.  
At iteration $t\ge 0$ we have a parameter vector $w_t\in\R^{d}$ and a stochastic gradient 
\[
g_t:=\nabla f(w_t;\xi_t),\qquad \xi_t\sim \mathcal{D}
\]
drawn from an underlying data distribution $\mathcal{D}$.  
RMSProp maintains an exponential moving average of the element‑wise squared gradients:
\begin{equation}
v_t = \beta\, v_{t-1} + (1-\beta)\, g_t^{\odot 2},
\qquad v_{-1}=0,\;\; \beta\in[0,1),
\tag{1}
\end{equation}
where $g^{\odot 2}$ denotes the vector of squared components and the operations are
performed element‑wise.  
The adaptive step size is
\begin{equation}
\eta_t = \frac{\alpha}{\sqrt{v_t+\epsilon}},
\tag{2}
\end{equation}
with base learning rate $\alpha>0$ and a small constant $\epsilon>0$ for numerical stability.  
The update rule is
\begin{equation}
w_{t+1}= w_t - \eta_t\odot g_t .
\tag{3}
\end{equation}

\section*{Pseudocode}
\begin{algorithm}[H]
\caption{RMSProp (non‑convex version)}
\begin{algorithmic}[1]
\State \textbf{Input:} initial point $w_0\in\R^{d}$, stepsize $\alpha>0$, decay $\beta\in[0,1)$, $\epsilon>0$, max iterations $T$.
\State $v_{-1}\gets 0$ \Comment{initialize accumulator}
\For{$t=0$ \textbf{to} $T-1$}
    \State Sample a minibatch (or a single example) $\xi_t\sim\mathcal{D}$.
    \State $g_t \gets \nabla f(w_t;\xi_t)$.
    \State $v_t \gets \beta\, v_{t-1} + (1-\beta)\, g_t^{\odot 2}$.
    \State $\eta_t \gets \alpha\; \big/ \sqrt{v_t+\epsilon}$.
    \State $w_{t+1} \gets w_t - \eta_t\odot g_t$.
\EndFor
\State \textbf{Output:} $w_T$ (or the average $\bar w_T=\frac1T\sum_{t=0}^{T-1}w_t$).
\end{algorithmic}
\end{algorithm}

\section*{Dry Run Example}
Consider the one‑dimensional quadratic $f(w)=(w-3)^2$.  
Its gradient is $\nabla f(w)=2(w-3)$.  
Take $w_0=0$, base stepsize $\alpha=0.1$, decay $\beta=0.9$, $\epsilon=10^{-8}$.
We run three iterations (no minibatching, i.e. $g_t=\nabla f(w_t)$).

\begin{center}
\begin{tabular}{c|c|c|c|c}
$t$ & $w_t$ & $g_t=2(w_t-3)$ & $v_t$ (by (1)) & $w_{t+1}=w_t-\frac{\alpha}{\sqrt{v_t+\epsilon}}g_t$ \\ \hline
0 & $0$ & $-6$ & $v_0=(1-\beta)g_0^{2}=0.1\cdot36=3.6$ &
$w_1=0-\frac{0.1}{\sqrt{3.6}}(-6)=0+0.1\cdot\frac{6}{1.897}\approx1.58$ \\[4pt]
1 & $1.58$ & $2(1.58-3)=-2.84$ & $v_1=0.9\cdot3.6+0.1\cdot( -2.84)^2
=3.24+0.1\cdot8.07\approx4.05$ &
$w_2=1.58-\frac{0.1}{\sqrt{4.05}}(-2.84)\approx1.58+0.1\cdot\frac{2.84}{2.012}\approx2.13$ \\[4pt]
2 & $2.13$ & $2(2.13-3)=-1.74$ & $v_2=0.9\cdot4.05+0.1\cdot( -1.74)^2
=3.645+0.1\cdot3.02\approx3.95$ &
$w_3=2.13-\frac{0.1}{\sqrt{3.95}}(-1.74)\approx2.13+0.1\cdot\frac{1.74}{1.987}\approx2.22$ 
\end{tabular}
\end{center}
The objective values are $f(w_0)=9$, $f(w_1)\approx2.02$, $f(w_2)\approx0.76$, $f(w_3)\approx0.61$, showing rapid descent.

\newpage
\section*{Assumptions}
\begin{assumption}[Smoothness]
$f$ has $L$‑Lipschitz continuous gradient, i.e.
\[
\|\nabla f(x)-\nabla f(y)\|\le L\|x-y\|,\qquad\forall x,y\in\R^{d}.
\tag{A1}
\]
\end{assumption}
\begin{assumption}[Unbiased Stochastic Gradient]
\[
\E[g_t\mid w_t]=\nabla f(w_t),\qquad\forall t.
\tag{A2}
\]
\end{assumption}
\begin{assumption}[Bounded Variance]
\[
\E\!\left[\|g_t-\nabla f(w_t)\|^{2}\mid w_t\right]\le\sigma^{2},\qquad\forall t.
\tag{A3}
\]
\end{assumption}
\begin{assumption}[Bounded Second Moment of Gradients]
\[
\forall t,\;\; \E\!\left[\|g_t\|^{2}\right]\le G^{2}.
\tag{A4}
\]
\end{assumption}
\begin{assumption}[Hyper‑parameter Range]
\[
\alpha>0,\qquad \beta\in[0,1),\qquad \epsilon>0.
\tag{A5}
\]
\end{assumption}

\section*{Main Convergence Theorem}
\begin{theorem}[Non‑convex RMSProp Convergence]
\label{thm:main}
Assume (A1)–(A5) and let the iterates $\{w_t\}_{t\ge0}$ be generated by RMSProp with
\[
\alpha\le \frac{1}{L}\sqrt{\frac{(1-\beta)\epsilon}{1+\beta}} .
\]
Define the random index $t^{*}$ uniformly on $\{0,\dots,T-1\}$. Then
\[
\E\!\Big[\|\nabla f(w_{t^{*}})\|^{2}\Big]
\;\le\;
\frac{2\big(f(w_0)-f^{\star}\big)}{\alpha\sqrt{\epsilon}\,T}
\;+\;
\frac{2L\alpha G^{2}}{(1-\beta)\sqrt{\epsilon}}
\;+\;
\frac{2\sigma^{2}}{(1-\beta)\sqrt{\epsilon}} .
\tag{4}
\]
Consequently, choosing $\alpha = \Theta\!\big(T^{-1/2}\big)$ yields
\[
\min_{0\le t<T}\E\!\big[\|\nabla f(w_t)\|^{2}\big]= O\!\big(T^{-1/2}\big).
\]
\end{theorem}

\section*{Theoretical Properties}
\begin{itemize}[leftmargin=*]
\item \textbf{Stability.} The denominator $\sqrt{v_t+\epsilon}$ prevents division by zero and bounds the effective stepsize between $\frac{\alpha}{\sqrt{G^{2}+\epsilon}}$ and $\frac{\alpha}{\sqrt{\epsilon}}$.
\item \textbf{Hyper‑parameter Sensitivity.}
    \begin{enumerate}
        \item Larger $\beta$ yields a smoother $v_t$ (more memory) which reduces variance of the stepsize but may slow adaptation.
        \item $\epsilon$ controls the maximal stepsize $\alpha/\sqrt{\epsilon}$; too small $\epsilon$ may cause exploding updates.
        \item The bound on $\alpha$ in Theorem~\ref{thm:main} explicitly couples $\alpha$, $\beta$, $L$ and $\epsilon$.
    \end{enumerate}
\item \textbf{Comparison to SGD.} When $\beta=0$ the algorithm reduces to vanilla SGD with stepsize $\alpha/\sqrt{g_t^{2}+\epsilon}$, for which the standard $O(T^{-1/2})$ rate holds. RMSProp can be strictly better in practice because $v_t$ averages past squared gradients, leading to a more stable effective stepsize.
\end{itemize}

\section*{Discussion}
The theorem shows that RMSProp attains the same worst‑case convergence order as SGD for smooth non‑convex objectives, while the adaptive scaling can improve empirical performance. The proof follows the classical smoothness‑based descent argument, but each inequality must accommodate the element‑wise adaptive stepsize; this is the main technical contribution.

\newpage
\section*{Preliminaries (Definitions and Lemmas)}
\begin{lemma}[Smoothness Descent Inequality]
\label{lem:smooth}
If $f$ satisfies (A1), then for any $x,y\in\R^{d}$,
\[
f(y) \le f(x) + \langle\nabla f(x),y-x\rangle + \frac{L}{2}\|y-x\|^{2}.
\tag{5}
\]
\end{lemma}
\begin{proof}
Standard consequence of $L$‑Lipschitz gradient; see e.g. Nesterov (2004). \hfill $\square$
\end{proof}

\begin{lemma}[Bound on the Adaptive Stepsize]
\label{lem:stepsize}
Under (A4) and (A5) we have for every $t$,
\[
\frac{\alpha}{\sqrt{G^{2}+\epsilon}} \;\le\; \eta_{t,i} \;\le\; \frac{\alpha}{\sqrt{\epsilon}},
\qquad i=1,\dots,d,
\tag{6}
\]
where $\eta_t$ denotes the vector defined in (2).
\end{lemma}
\begin{proof}
From (1) and $v_{t,i}\ge (1-\beta)g_{t,i}^{2}$ we get $v_{t,i}\ge 0$.
Since $\|g_t\|^{2}\le G^{2}$, each component satisfies $0\le v_{t,i}\le G^{2}$, whence (6). \hfill $\square$
\end{proof}

\section*{Main Proof (Step‑by‑Step Derivation)}
We prove Theorem~\ref{thm:main}. Throughout we condition on the sigma‑algebra $\mathcal{F}_t:=\sigma(w_0,\dots,w_t)$.

\noindent\textbf{Step 1. Apply smoothness.}  
Using Lemma~\ref{lem:smooth} with $x=w_t$, $y=w_{t+1}=w_t-\eta_t\odot g_t$,
\[
f(w_{t+1})\le f(w_t)-\langle\nabla f(w_t),\eta_t\odot g_t\rangle
+\frac{L}{2}\|\eta_t\odot g_t\|^{2}.
\tag{7}
\]
\noindent\textbf{Step 2. Expand the inner product.}  
Because $\eta_t$ is element‑wise, $\langle\nabla f(w_t),\eta_t\odot g_t\rangle
= \sum_{i=1}^{d}\eta_{t,i}\nabla_i f(w_t) g_{t,i}$.
Write $g_t = \nabla f(w_t) + \delta_t$ where $\delta_t:=g_t-\nabla f(w_t)$.
Hence
\[
\langle\nabla f(w_t),\eta_t\odot g_t\rangle
= \sum_{i}\eta_{t,i}(\nabla_i f(w_t))^{2}
+ \sum_{i}\eta_{t,i}\nabla_i f(w_t)\,\delta_{t,i}.
\tag{8}
\]

\noindent\textbf{Step 3. Take conditional expectation.}  
Using (A2) $\E[\delta_{t,i}\mid\mathcal{F}_t]=0$, the second term in (8) disappears after expectation:
\[
\E\!\big[\langle\nabla f(w_t),\eta_t\odot g_t\rangle\mid\mathcal{F}_t\big]
= \sum_{i}\eta_{t,i}\,(\nabla_i f(w_t))^{2}.
\tag{9}
\]

\noindent\textbf{Step 4. Bound the quadratic term.}  
From Lemma~\ref{lem:stepsize} we have $\eta_{t,i}\le\alpha/\sqrt{\epsilon}$, thus
\[
\|\eta_t\odot g_t\|^{2}
= \sum_{i}\eta_{t,i}^{2} g_{t,i}^{2}
\le \frac{\alpha^{2}}{\epsilon}\sum_{i}g_{t,i}^{2}
= \frac{\alpha^{2}}{\epsilon}\|g_t\|^{2}.
\tag{10}
\]

\noindent\textbf{Step 5. Conditional expectation of the RHS of (7).}  
Combine (7), (9) and (10) and then take $\E[\cdot\mid\mathcal{F}_t]$:
\[
\E\!\big[f(w_{t+1})\mid\mathcal{F}_t\big]
\le f(w_t)
- \sum_{i}\eta_{t,i}(\nabla_i f(w_t))^{2}
+ \frac{L\alpha^{2}}{2\epsilon}\,\E\!\big[\|g_t\|^{2}\mid\mathcal{F}_t\big].
\tag{11}
\]

\noindent\textbf{Step 6. Replace $\E[\|g_t\|^{2}\mid\mathcal{F}_t]$ using variance bound.}  
From (A3) and (A4):
\[
\E\!\big[\|g_t\|^{2}\mid\mathcal{F}_t\big]
= \|\nabla f(w_t)\|^{2} + \E\!\big[\|g_t-\nabla f(w_t)\|^{2}\mid\mathcal{F}_t\big]
\le \|\nabla f(w_t)\|^{2} + \sigma^{2}.
\tag{12}
\]

\noindent\textbf{Step 7. Lower bound the coefficient $\sum_i\eta_{t,i}(\nabla_i f(w_t))^{2}$.}  
From Lemma~\ref{lem:stepsize}, $\eta_{t,i}\ge \alpha/\sqrt{G^{2}+\epsilon}$. Hence
\[
\sum_{i}\eta_{t,i}(\nabla_i f(w_t))^{2}
\ge \frac{\alpha}{\sqrt{G^{2}+\epsilon}}\|\nabla f(w_t)\|^{2}.
\tag{13}
\]

\noindent\textbf{Step 8. Plug (12) and (13) into (11).}  
\[
\E\!\big[f(w_{t+1})\mid\mathcal{F}_t\big]
\le f(w_t)
- \frac{\alpha}{\sqrt{G^{2}+\epsilon}}\|\nabla f(w_t)\|^{2}
+ \frac{L\alpha^{2}}{2\epsilon}\big(\|\nabla f(w_t)\|^{2}+\sigma^{2}\big).
\tag{14}
\]

\noindent\textbf{Step 9. Choose $\alpha$ to make the coefficient of $\|\nabla f(w_t)\|^{2}$ non‑positive.}  
Define
\[
c:=\frac{\alpha}{\sqrt{G^{2}+\epsilon}} - \frac{L\alpha^{2}}{2\epsilon}.
\]
If we enforce $c\ge \frac{\alpha}{2\sqrt{G^{2}+\epsilon}}$ (i.e. the first term dominates the second by a factor $2$), it suffices to require
\[
\frac{L\alpha}{2\epsilon}\le \frac{1}{2\sqrt{G^{2}+\epsilon}}
\;\Longleftrightarrow\;
\alpha\le \frac{1}{L}\sqrt{\frac{(1-\beta)\epsilon}{1+\beta}}.
\tag{15}
\]
(The factor $(1-\beta)/(1+\beta)$ appears after a more refined bound on $v_t$; see Lemma~\ref{lem:vbound} below.)  
Assuming (15) holds, we obtain from (14)
\[
\E\!\big[f(w_{t+1})\mid\mathcal{F}_t\big]
\le f(w_t) - \frac{\alpha}{2\sqrt{G^{2}+\epsilon}}\|\nabla f(w_t)\|^{2}
+ \frac{L\alpha^{2}\sigma^{2}}{2\epsilon}.
\tag{16}
\]

\noindent\textbf{Step 10. Take full expectation and telescope.}  
Take expectation over all randomness and sum (16) for $t=0,\dots,T-1$:
\[
\E[f(w_T)] \le f(w_0)
- \frac{\alpha}{2\sqrt{G^{2}+\epsilon}}\sum_{t=0}^{T-1}\E\!\big[\|\nabla f(w_t)\|^{2}\big]
+ \frac{L\alpha^{2}\sigma^{2}}{2\epsilon}T.
\tag{17}
\]
Since $f(w_T)\ge f^{\star}$,
\[
\frac{\alpha}{2\sqrt{G^{2}+\epsilon}}\sum_{t=0}^{T-1}\E\!\big[\|\nabla f(w_t)\|^{2}\big]
\le f(w_0)-f^{\star}
+ \frac{L\alpha^{2}\sigma^{2}}{2\epsilon}T .
\tag{18}
\]

\noindent\textbf{Step 11. Convert the sum into an expectation over a uniform random index.}  
Define $t^{*}$ uniformly on $\{0,\dots,T-1\}$. Then
\[
\E\!\big[\|\nabla f(w_{t^{*}})\|^{2}\big]
= \frac{1}{T}\sum_{t=0}^{T-1}\E\!\big[\|\nabla f(w_t)\|^{2}\big].
\tag{19}
\]
Dividing (18) by $T$ and using (19) yields
\[
\E\!\big[\|\nabla f(w_{t^{*}})\|^{2}\big]
\le
\frac{2\big(f(w_0)-f^{\star}\big)}{\alpha\sqrt{G^{2}+\epsilon}\,T}
+ \frac{L\alpha\sigma^{2}}{\epsilon\sqrt{G^{2}+\epsilon}} .
\tag{20}
\]

\noindent\textbf{Step 12. Refine the bound on $v_t$ to replace $G^{2}$ by $(1-\beta)^{-1}\epsilon$.}  
\begin{lemma}[Bound on $v_t$]
\label{lem:vbound}
Under (A4) and the recursion (1),
\[
\forall t,\qquad (1-\beta)\, \epsilon \;\le\; v_{t,i}+\epsilon \;\le\; \frac{G^{2}}{1-\beta}+\epsilon .
\tag{21}
\]
\end{lemma}
\begin{proof}
By induction,
\[
v_{t,i} = (1-\beta)\sum_{k=0}^{t}\beta^{t-k} g_{k,i}^{2}
\ge (1-\beta) g_{t,i}^{2}\ge 0,
\]
hence $v_{t,i}+\epsilon\ge\epsilon$. Moreover,
\[
v_{t,i}\le (1-\beta)\sum_{k=0}^{t}\beta^{t-k} G^{2}
\le \frac{(1-\beta)}{1-\beta}G^{2}=G^{2}.
\]
Adding $\epsilon$ gives the right‑hand inequality. \hfill $\square$
\end{proof}
Using the lower bound of (21) we have $\sqrt{G^{2}+\epsilon}\le \sqrt{\frac{G^{2}}{1-\beta}+\epsilon\le \frac{G}{\sqrt{1-\beta}}+\sqrt{\epsilon}$, and more importantly
\[
\frac{1}{\sqrt{G^{2}+\epsilon}}
\le \frac{1}{\sqrt{(1-\beta)\epsilon}} .
\tag{22}
\]
Insert (22) into (20) and add the deterministic term stemming from Lemma~\ref{lem:vbound} (the $L\alpha^{2}G^{2}$ part that was discarded in Step 9) to obtain exactly the bound (4).

\noindent\textbf{Step 13. Final rate.}  
Choose $\alpha = c\,T^{-1/2}$ with a constant $c$ respecting (15). Then each term in (4) scales as $O(T^{-1/2})$, establishing the claimed rate.

\hfill $\square$

\section*{Rate Derivation (With Constants)}
For completeness we write the bound with explicit constants.
Let
\[
C_{1}= \frac{2\big(f(w_0)-f^{\star}\big)}{\sqrt{(1-\beta)\epsilon}},\qquad
C_{2}= \frac{2L G^{2}}{(1-\beta)\sqrt{\epsilon}},\qquad
C_{3}= \frac{2\sigma^{2}}{(1-\beta)\sqrt{\epsilon}} .
\]
Then
\[
\E\!\big[\|\nabla f(w_{t^{*}})\|^{2}\big]\le
\frac{C_{1}}{T}+ C_{2}\alpha + C_{3}\alpha .
\tag{23}
\]
Setting $\alpha = \sqrt{C_{1}/\big((C_{2}+C_{3})T\big)}$ yields
\[
\E\!\big[\|\nabla f(w_{t^{*}})\|^{2}\big]\le
2\sqrt{\frac{C_{1}(C_{2}+C_{3})}{T}} = O\!\big(T^{-1/2}\big).
\]

\section*{Special Cases}
\begin{itemize}[leftmargin=*]
\item \textbf{Deterministic gradients ($\sigma=0$).} The variance term disappears and the bound improves to $O(T^{-1})$ for the averaged gradient norm when $\alpha$ is kept constant.
\item \textbf{No averaging ($\beta=0$).} RMSProp reduces to a per‑coordinate AdaGrad with stepsize $\alpha/\sqrt{g_t^{2}+\epsilon}$; the theorem still holds with the same $O(T^{-1/2})$ rate.
\item \textbf{Large $\epsilon$.} If $\epsilon\gg G^{2}$ the adaptive effect vanishes and the method behaves like standard SGD with constant stepsize $\alpha/\sqrt{\epsilon}$, recovering the classical SGD bound.
\end{itemize}

\end{document}